\chapter{Architectural Landscape of SVG Generation}
\label{ch:Landscape}

SVG generation is not a single problem with a single solution. 
Some methods treat vector graphics as optimization targets and 
adjust geometric parameters until a rendered image matches a 
prompt or a reference. Others learn from datasets and generate 
SVGs through a trained decoder. A smaller group avoids explicit 
geometric primitives entirely and uses neural functions to 
define shapes indirectly. These differences are not just 
technical details --- they reflect fundamentally different 
ideas about what a good SVG output should look like, and they 
lead to very different trade-offs between visual quality, 
structural organization, and editability.

To make these trade-offs visible, this thesis analyzes the 
selected models across three dimensions: how they represent 
vector graphics internally, how they supervise the generation 
process, and what the results look like in terms of structure 
and usability. This chapter introduces the thirteen models 
that form the basis of that analysis, organized into five 
families according to their shared design assumptions. Each 
family represents a distinct way of thinking about SVG 
generation, and understanding the differences between them 
is what makes the later analysis possible.

The five families cover optimization-based vectorization, 
diffusion-guided generation, neural implicit representations, 
dataset-trained generative models, and neural vectorization 
systems. The boundaries between them are not always sharp --- 
some recent models deliberately combine ideas from more than 
one family. But the grouping is still useful because it keeps 
the focus on what actually differs between approaches, rather 
than treating each model as a separate contribution to be 
described in isolation.



Methods that generate SVG graphics can be grouped into several 
research directions, each addressing the structural challenges of 
vector representations differently.

A first line of work treats vector drawings as ordered sequences 
of commands or strokes. Inspired by language modelling, these 
methods learn to predict drawing operations step by step using 
recurrent or transformer-based architectures. SketchRNN~\cite{haDrawTogetherNeural2017} 
showed that neural networks can learn stroke dynamics and produce 
plausible vector-style sketches. Later models extended this to SVG 
command prediction, including DeepSVG~\cite{carlierDeepSVGHierarchicalGenerative2020}, 
Sketchformer~\cite{ribeiroSketchformerTransformerbasedRepresentation2020}, 
and subsequent sequence-based works~\cite{ganinSpiralGeneratingCurved2018, 
lopesLearningRepresentationsSketch2019}. However, long command 
chains make training difficult, and keeping global visual coherence 
and structural correctness at the same time becomes challenging. 
SVG datasets also tend to cover simple graphics --- icons, logos, 
basic illustrations. There are no large-scale datasets for complex 
scenes, which limits how far these models can generalize.

A second direction moves the learning process into the image domain 
through differentiable rasterization. Here, a set of vector 
parameters is initialized and then iteratively optimized so that 
the rendered image matches a target raster image or a textual 
description. DiffVG~\cite{liDifferentiableVectorGraphics2020a} 
makes this possible by allowing gradients to flow back from pixel 
space to vector parameters. Models such as 
LIVE~\cite{maLayerwiseImageVectorization2022} iteratively refine 
layered Bézier paths to approximate images, while 
VectorFusion~\cite{jainVectorFusionTexttoSVGSymbol2023} and 
SVGDreamer~\cite{xingSVGDreamerTextGuided2024} combine 
differentiable rasterization with diffusion-based supervision to 
enable text-to-vector generation. These approaches can produce 
visually strong results, but they may also produce unstable 
geometry and messy paths, which makes editing difficult.

A third direction avoids directly optimizing discrete primitives. 
Instead of updating Bézier control points, paths, or shapes, 
these methods model vector graphics through neural networks and 
run optimization in neural parameter space. Frameworks such as 
NeuralSVG~\cite{vinkerNeuralSVGInstancededbasedSVG2025} and 
NiVeL~\cite{thamizharasanNIVeLNeuralImplicit2024} represent shape 
boundaries through a neural function --- typically an MLP --- 
whose weights are optimized using image-space supervision. This 
can be thought of as a hybrid approach that bridges vector 
optimization and neural representation learning.

More recent work has shifted toward large language models (LLMs) 
and multimodal large language models (MLLMs). Since SVG is 
fundamentally a structured text format, language models are a 
natural fit for this task. Methods such as 
StarVector~\cite{rodriguezStarVectorGeneratingScalable2025}, 
OmniSVG~\cite{yangOmniSVGUnifiedScalable2025}, 
LLM4SVG~\cite{xingEmpoweringLLMsUnderstand2025}, and 
InternSVG~\cite{wangInternSVGUnifiedSVG2025} treat SVG generation 
as code synthesis, using transformer architectures pretrained on 
large-scale text and image data to generate XML commands from 
textual or visual inputs. Compared to earlier sequence-based 
models, these approaches show stronger alignment between text 
prompts and visual outputs and cover a wider range of tasks 
including understanding, generation, and editing.

However, this paradigm has a problem that is directly relevant to 
this thesis. SVG files are long --- even simple graphics produce 
hundreds of lines of XML --- and transformer models struggle to 
maintain structural consistency across sequences of that length. 
More fundamentally, generation is distributed across large neural 
architectures rather than defined through explicit vector 
operations, which makes it difficult to trace how geometric 
structure emerges from architectural decisions. For these reasons, 
LLM-based SVG generation is treated here as an emerging direction 
rather than a primary subject of analysis. The following chapters 
focus on optimization-based, neural, and hybrid approaches, where 
the relationship between representation and structural outcomes 
can be examined more directly.

\section{Optimization-Based SVG Generation}

One of the most influential paradigms in modern SVG generation treats 
vector graphics not as outputs predicted by a model, but as variables 
to be optimized. In this formulation, the SVG itself becomes the 
parameter space of the problem. Instead of generating vector primitives 
directly, the system initializes a set of shapes and iteratively refines 
their parameters so that the rendered result matches a target image or 
structural signal.

This approach became practical with the introduction of differentiable 
rasterization. As introduced in Chapter~\ref{ch:Background}, DiffVG~\cite{liDifferentiableVectorGraphics2020a} 
provides a bridge between vector parameter space and pixel space by 
making the rasterization step differentiable. In the optimization-based 
paradigm, this means that a loss defined on the rendered image --- 
whether a pixel reconstruction loss or a semantic similarity score --- 
can propagate gradients directly back to vector parameters such as 
Bézier control points, stroke widths, and colors. Without this 
capability, iterative refinement of vector primitives using 
gradient-based methods would not be possible. DiffVG therefore 
functions not as a generation model in itself, but as the foundational 
rendering mechanism on which the entire optimization-based family 
depends.

Methods in this family share several defining characteristics.
Supervision is typically defined through reconstruction or structural 
objectives rather than dataset-level learning.  
They rely on explicit parametric representations of vector 
graphics, meaning that primitives are theoretically accessible for 
editing throughout the generation process. In practice, however, 
the number of paths produced by optimization-based methods often 
makes meaningful editing difficult, as primitives tend to become 
entangled or fragmented when supervision operates primarily in 
raster space~\cite{xingSVGDreamer2024, zhuSAMVGMultistageImage2023}.

LIVE~\cite{maLayerwiseImageVectorization2022} represents a canonical example 
of this paradigm. The method introduces a layer-wise optimization 
strategy in which vector primitives are refined progressively, allowing 
structure to emerge gradually instead of being predicted at once. This 
coarse-to-fine refinement encourages hierarchical organization and often 
leads to SVG outputs that remain interpretable and editable. However, 
because optimization focuses primarily on visual reconstruction, the 
resulting graphics may contain redundant paths or fragmented geometry 
when structural constraints are not explicitly enforced.

SAMVG~\cite{zhuSAMVGMultistageImage2023} extends this optimization 
formulation by introducing semantic initialization through segmentation 
masks. As illustrated in \includegraphics[width=0.85\textwidth]{LIVEnSAM.png}, rather than 
starting from random primitives, the method leverages the Segment 
Anything Model to guide the initial placement and grouping of vector 
elements. Both LIVE and SAMVG share the same core optimization loop --- 
initialized SVG parameters are passed through DiffVG, rendered to an 
image, and updated via pixel reconstruction loss against a target. The 
critical difference lies in initialization: where LIVE begins from 
randomly placed paths, SAMVG uses semantically meaningful regions 
derived from segmentation. This highlights a central insight within this 
family: initialization strategy strongly influences structural outcomes, 
independent of the optimization mechanism itself.

From an analytical perspective, optimization-based methods occupy an 
important position in the SVG generation landscape. They provide the 
strongest connection between vector parameters and final structure 
because primitives remain explicit throughout the pipeline. This makes 
them particularly relevant when editability is a primary objective. At 
the same time, the reliance on iterative refinement introduces practical 
limitations related to inference time, stability, and scalability to 
complex scenes.

More broadly, this family establishes a baseline for understanding 
subsequent approaches. Many later methods retain the optimization 
formulation while modifying supervision signals, representation spaces, 
or initialization strategies. As a result, optimization-based vector 
generation serves as a reference point for analyzing how additional 
architectural components influence the balance between visual realism 
and structural quality.


\section{Diffusion-Guided Optimization-based SVG Generation}

While classical optimization-based methods rely primarily on reconstruction objectives, more recent approaches extend this paradigm by introducing semantic supervision from pretrained diffusion models. Instead of matching a fixed target image, vector parameters are optimized so that the rendered SVG becomes consistent with the distribution learned by a text-conditioned generative model. This shift fundamentally changes the role of supervision in SVG generation.

The key mechanism enabling this transition is Score Distillation Sampling (SDS), which allows gradients derived from diffusion models to guide optimization in vector parameter space. In practice, the SVG is rendered into a raster image, evaluated by a diffusion model conditioned on a text prompt, and then updated using gradients that encourage semantic alignment. As a result, vector graphics can be generated from textual descriptions without requiring paired SVG datasets.

Methods such as VectorFusion, SVGDreamer, and SVGDreamer++ exemplify this family. These approaches retain the explicit parametric representation of vector graphics introduced by classical optimization methods but replace reconstruction losses with semantic priors. This combination enables the generation of visually rich and diverse SVG outputs that reflect the expressive capabilities of modern diffusion models.

However, this improvement in visual realism introduces new structural challenges. Because diffusion models operate in raster space, the gradients they provide do not explicitly enforce geometric consistency or hierarchy. Optimization therefore prioritizes appearance, often leading to redundant primitives, unstable paths, or overly dense representations. This reveals a central trade-off in diffusion-guided vector optimization: stronger semantic guidance improves realism but may degrade structural clarity and editability.

Different models within this family attempt to address this limitation through architectural refinements. SVGDreamer introduces strategies for primitive control and structured optimization, while SVGDreamer++ further explores hierarchical refinement mechanisms that adapt the number and organization of vector elements during generation. These efforts reflect an emerging recognition that semantic supervision alone is insufficient for producing high-quality vector structure and must be complemented by mechanisms that regulate geometry.

From a broader perspective, diffusion-guided optimization represents a transitional paradigm in SVG generation. It demonstrates that strong generative priors can be transferred from raster models into vector domains, dramatically improving visual expressiveness. At the same time, it exposes the limitations of supervision that is detached from vector structure. Understanding this tension is essential for analyzing why some methods achieve photorealistic SVG outputs while still struggling to produce representations that remain easily editable.

This family therefore plays a critical role in the architectural landscape. It highlights how changes in supervision can reshape optimization behavior and directly influence structural outcomes, setting the stage for approaches that move beyond explicit parameter optimization toward alternative representation strategies.


\section{Neural Implicit and Hybrid SVG Generation}

An alternative direction in SVG generation moves away from directly optimizing explicit vector primitives and instead represents graphics through neural functions. In this formulation, geometry is not stored as a fixed set of paths but is defined implicitly by the parameters of a neural network. Vector structure therefore emerges from evaluating a learned function rather than being explicitly specified from the beginning.

Methods such as NeuralSVG, NiVeL, and T2V-NPR exemplify this paradigm. These approaches combine optimization with neural parameterization, where the weights of a neural network — typically a multilayer perceptron — define shape boundaries, layers, or path distributions. Instead of adjusting Bézier control points directly, optimization occurs in neural parameter space, and the resulting function is interpreted as a vector graphic.

This shift introduces an important conceptual difference compared to classical optimization-based methods. While parametric approaches treat primitives as independent elements, neural implicit representations encourage global coherence because all geometry is generated through a shared function. As a result, structure may emerge in a more stable and organized manner, often producing layered outputs that are easier to interpret and manipulate.

NeuralSVG illustrates this idea by encoding scene structure within neural weights that are optimized using semantic guidance. NiVeL further explores decomposable layered representations, allowing graphics to be constructed as combinations of implicit components. T2V-NPR introduces neural path representations that operate in a learned latent space, providing additional constraints that improve optimization stability and path validity.

From the perspective of structural quality, this family highlights a key insight: representation itself can act as a structural prior. When geometry is produced through a continuous neural function, the model implicitly restricts how primitives can be arranged, which may reduce fragmentation and encourage consistent layering. This contrasts with purely parametric optimization, where primitives can drift independently.

At the same time, this approach introduces new trade-offs. Because optimization occurs in neural parameter space, inference typically remains iterative and may require substantial computation to converge. The computational cost is therefore influenced not only by representation but also by supervision signals, optimization strategy, and rendering complexity. In practice, these methods often balance improved structural coherence with increased inference time.

More broadly, neural implicit and hybrid approaches represent an attempt to bridge the gap between explicit vector optimization and learned generative structure. They suggest that achieving editable SVG outputs may depend not only on supervision strength but also on how structure is encoded within the representation itself. This observation plays an important role in the subsequent analysis of representation–supervision interactions across SVG generation methods.



\section{Dataset-Trained SVG Generation}

In contrast to optimization-based approaches, another line of research treats SVG generation as a learning problem in which structure is acquired from data rather than constructed through iterative refinement. In this paradigm, models are trained on collections of vector graphics to learn statistical regularities in geometry, composition, and drawing structure. Generation then occurs through a forward decoding process, where a learned representation is transformed directly into SVG primitives.

Models such as DeepSVG and DeepIcon illustrate this dataset-driven formulation. These approaches learn latent spaces over vector graphics, allowing SVG commands or path sequences to be generated through autoregressive decoding or hierarchical prediction. Because structural patterns are learned from data, the resulting outputs often exhibit consistent geometry and stylistic coherence without requiring explicit optimization during inference.

SVGFusion extends this idea by combining latent representation learning with diffusion modeling. Instead of optimizing vector parameters per instance, the model first learns a continuous latent space of SVG graphics from a dataset and then applies a diffusion transformer to generate new samples in that space. SVG generation therefore becomes a process of sampling and decoding rather than iterative parameter refinement. This distinguishes SVGFusion from diffusion-guided optimization methods, which rely on semantic gradients rather than learned vector distributions.

From a structural perspective, dataset-trained methods offer a different set of advantages. Because the model learns common drawing patterns, generated SVGs may exhibit more regular organization and fewer redundant primitives compared to unconstrained optimization approaches. Latent representations can also encourage smooth interpolation and variation, making these methods well suited for icon generation and style-consistent vector synthesis.

However, this reliance on training data introduces its own limitations. Structural quality is bounded by the diversity and characteristics of the dataset, which may restrict the ability to generalize to complex scenes or highly detailed compositions. Furthermore, although inference is typically faster than iterative optimization, forward generation reduces direct control over individual primitives, which can affect editability.

This family therefore highlights an alternative route toward structured SVG generation: rather than enforcing structure through optimization constraints, models attempt to internalize it through learning. Understanding this distinction is essential for analyzing how representation and supervision jointly influence the balance between consistency, realism, and user-level control.



\section{Neural Vectorization and Generative Systems}

A related but distinct direction in SVG research focuses on vectorization, where the objective is to translate raster images or multimodal inputs into structured vector representations. Unlike generative models that aim to sample new graphics, vectorization approaches emphasize reconstructing existing visual content while preserving geometric structure. This shifts the problem from distribution modeling toward representation inference.

Im2Vec represents one of the earlier neural formulations of this task, where models learn to predict parametric vector primitives directly from images using reconstruction-based supervision. In this setting, vector graphics are treated as outputs of an encoder–decoder pipeline, and generation occurs through a single forward pass. This approach enables efficient inference and supports large-scale vectorization, but it also introduces challenges related to structural ambiguity, since multiple vector representations may correspond to the same raster image.

More recent systems such as LayerTracer extend this idea by modeling the construction process of vector graphics rather than predicting primitives independently. Instead of producing a flat set of paths, the model learns layered representations that approximate how designers build graphics step by step. Diffusion-based architectures are used to generate intermediate structures that guide vector formation, allowing the system to capture relationships between elements and improve hierarchical organization.

This distinction between primitive prediction and construction modeling is particularly relevant for editability. When vectorization is framed as a direct regression problem, outputs may approximate appearance effectively but lack meaningful grouping. By contrast, approaches that model layered structure attempt to embed designer logic into the generation process, potentially producing SVG outputs that are easier to modify and reuse.

From a broader perspective, neural vectorization systems highlight the importance of task formulation. While generative models prioritize diversity and semantic alignment, vectorization approaches prioritize fidelity and structural reconstruction. This often leads to different trade-offs in representation choice, supervision design, and inference behavior. Forward prediction enables faster generation compared to iterative optimization, but structural guarantees depend heavily on how the model encodes relationships between elements.

This family therefore acts as a bridge between traditional algorithmic vectorization and modern generative pipelines. It demonstrates how learning-based methods can incorporate procedural knowledge and layered reasoning, offering insights into how designer-oriented SVG structure might be modeled directly within neural architectures.


\section{Synthesis and Transition to Analytical Dimensions}

The model families introduced in this chapter illustrate that SVG generation is not a single unified problem but a collection of design strategies that differ in how structure is represented, how supervision is defined, and how geometry is constructed. Each family highlights a distinct perspective on vector graphics generation, revealing implicit assumptions about what constitutes a successful SVG output.

Optimization-based methods emphasize direct control over geometry by treating vector parameters as optimization variables, often producing editable but computationally expensive representations. Diffusion-guided approaches demonstrate how strong semantic priors can dramatically improve visual realism, while simultaneously exposing the limitations of supervision that is detached from vector structure. Neural implicit methods shift structural decisions into representation space, suggesting that how geometry is encoded can influence stability and layering. Dataset-trained generative models instead learn structural regularities from data, trading explicit control for consistency and efficiency. Finally, neural vectorization systems illustrate how modeling the construction process itself can affect hierarchical organization and designer usability.

Taken together, these families reveal that structural properties of generated SVGs do not emerge from a single architectural component. Instead, they arise from the interaction between representation, supervision, and inference strategy. Similar visual outcomes may therefore be produced through fundamentally different pipelines, while small changes in supervision or representation can lead to significant differences in editability, redundancy, and computational cost.

This observation motivates the analytical perspective adopted in the following chapters. Rather than comparing models solely at the level of outputs, the remainder of this thesis examines the architectural dimensions that shape SVG structure. In particular, the analysis focuses on internal representation, supervision signals, optimization regimes, and their combined influence on editability and structural coherence.

By shifting the discussion from model families to architectural dimensions, the thesis aims to identify recurring design patterns and clarify which combinations of representation and supervision are most likely to produce SVG outputs that are both visually convincing and practically usable. This transition establishes the foundation for a comparative analysis that moves beyond individual methods toward a deeper understanding of how SVG generation systems are designed